\section{Related Work}
\label{sec:related-work}

\subsection{From statistical synthesis to deep generators}

Synthetic microdata predate deep learning. Multiple imputation and probabilistic graphical models generate records from estimated conditional distributions and remain attractive because their assumptions can be inspected \citep{raghunathan2003}. Copulas provide efficient dependence modelling but can struggle when discrete, continuous, multimodal, and rare-event variables interact nonlinearly \citep{wang2021}.

Generative adversarial networks (GANs) formulate generation as an adversarial game between a generator and a discriminator \citep{goodfellow2014}. MedGAN adapted this idea to high-dimensional discrete patient records using an autoencoder \citep{choi2017}. CTGAN then addressed mixed-type tables through mode-specific normalisation, conditional vectors, and training-by-sampling \citep{xu2019}. These mechanisms are directly relevant to health tables with skewed continuous variables and imbalanced categories, which is why we adopt CTGAN as the common generator throughout the experiments reported here.

Variational autoencoders (VAEs) optimise a likelihood-related lower bound and offer an explicit latent representation \citep{kingma2014}; TVAE was evaluated alongside CTGAN in the same tabular-synthesis framework \citep{xu2019}. More recently, diffusion models iteratively denoise samples, and TabDDPM has shown that separate diffusion treatments for numerical and categorical attributes can outperform several GAN and VAE baselines on general tabular benchmarks \citep{kotelnikov2023}. These results motivate diffusion models as baselines for future comparison, but they do not establish superiority for the small, imbalanced clinical datasets considered here, and no diffusion baseline is included in the present experiments (Section~\ref{sec:limitations}).

\subsection{Health-data synthesis and evaluation}

Health-synthesis studies repeatedly conclude that no single metric establishes usefulness or safety, and that evaluation should jointly cover resemblance, downstream utility, and privacy \citep{goncalves2020,hernandez2023,pezoulas2024}. Recent large-scale comparisons likewise report dataset- and metric-dependent rankings and a persistent fidelity/utility--privacy trade-off \citep{hernandez2025,nanevski2026}, and a systematic review of structured synthetic EHR generation documents continuing difficulties with longitudinal structure, rare events, reproducibility, and privacy evaluation \citep{ghosheh2024}.

Within this evaluation literature, univariate fidelity measures whether each synthetic column resembles its real counterpart, pairwise fidelity assesses correlation or contingency structure, and multivariate distinguishability tests whether a classifier can separate real from synthetic records. None of these is sufficient on its own for a downstream clinical task: a dataset can score well in aggregate because majority-class records dominate its statistics while the minority class is poorly represented, and conversely, aggressive selection can improve classification simply by producing easy prototypes rather than a faithful population.

Train-on-synthetic, test-on-real (\tstr{}) evaluation probes directly whether relations learned from synthetic data transfer to real observations, and is a stronger test than evaluating on synthetic data alone, because an internally self-consistent synthetic distribution can still be systematically wrong. \tstr{} is not, however, a privacy test, and it does not establish that synthetic-data estimates are inferentially unbiased. Privacy must be evaluated separately, through nearest-record, membership-inference, or attribute-inference analyses, or through a formal differential-privacy guarantee \citep{yale2019,dwork2006}: synthetic data are not anonymous by construction, and a high-capacity generator can memorise training records.

\subsection{Class imbalance and rare clinical outcomes}

Class imbalance changes both training and evaluation. Accuracy can remain high even when minority recall is poor; Macro-F1 weights classes equally, and the area under the precision--recall curve (AUPRC) is often more informative than the area under the ROC curve (AUROC) when the positive outcome is rare \citep{saito2015}. Conventional resampling methods, including random oversampling and SMOTE \citep{chawla2002}, can help, but interpolated or duplicated points may distort heterogeneous minority subgroups; broader reviews of imbalanced learning emphasise that sampling strategy, cost weighting, decision thresholds, and evaluation metric must be matched to the application \citep{he2009}.

Synthetic augmentation can increase sample count without increasing information. The relevant question is therefore not whether the generated dataset is balanced, but whether the minority conditional distribution $p(x\mid y=1)$ and its overlap with $p(x\mid y=0)$ are preserved well enough to improve performance on held-out real cases. \sepa{} targets this gap directly: it selects label-plausible candidates under an explicit separability criterion while retaining a realism term and the real target-class counts, so that class ratio is held constant across compared conditions.

\subsection{Causal structure and generative modelling}

Associational fidelity does not imply causal fidelity. A directed acyclic graph (DAG) encodes qualitative assumptions about data-generating relations, while causal-effect estimation additionally requires identification assumptions that cannot be learned from observational data alone \citep{pearl2009,spirtes2000}. Continuous-optimisation methods such as NOTEARS make DAG discovery scalable \citep{zheng2018}, but discovered graphs can be unstable under small samples, mixed variable types, measurement error, and hidden confounding -- all of which apply to the datasets used here.

CausalGAN showed that a known causal graph can organise adversarial generation \citep{kocaoglu2018}, and causal-effect VAEs illustrate how generative latent-variable models can encode causal assumptions \citep{louizos2017}. These methods motivate causally constrained synthesis in principle, but domain validation remains essential in practice. Accordingly, the \dacpf{} experiment reported in Section~\ref{sec:results} is described throughout as \emph{causal plausibility filtering}, not causal discovery or effect identification: it rewards consistency with predeclared, clinician-motivated anchor patterns, but it does not estimate or test a causal effect.

\subsection{Privacy--utility tension}

Differential privacy bounds the effect of any single training record on an algorithm's output \citep{dwork2006}. Within \emph{Computers in Biology and Medicine} specifically, LDP-GAN combines local differential privacy with medical-record synthesis and evaluates both utility and privacy \citep{kim2024ldpgan}, and a recent oncology framework in the same journal evaluates and selects among synthetic datasets for survival prediction across several quality dimensions \citep{christoforou2025}. These studies indicate the journal's interest in medically motivated synthesis accompanied by rigorous downstream and privacy evaluation, rather than in generator novelty alone -- a bar that the present work aims toward but, as discussed in Section~\ref{sec:limitations}, does not yet fully clear on the privacy dimension.

\subsection{Research gap}

Taken together, the literature leaves a practical gap at the intersection of four concerns: rare clinical outcomes, transparent post-generation control, causal plausibility, and cross-dataset validation. Existing generators optimise distribution learning; imbalance methods optimise prediction; causally informed generators typically assume a trusted graph; and evaluation frameworks diagnose synthetic outputs after generation rather than steering it. This work asks whether these components can be connected into an auditable pipeline in which each selection pressure has a measurable, attributable effect, and in which any improvement in \tstr{} utility is required to survive separate tests of fidelity, subgroup structure, causal consistency, external validity, and privacy.
